\IfFileExists{stacks-project.cls}{%
\documentclass{stacks-project}
}{%
\documentclass{amsart}
}

% For dealing with references we use the comment environment
\usepackage{verbatim}
\newenvironment{reference}{\comment}{\endcomment}
%\newenvironment{reference}{}{}
\newenvironment{slogan}{\comment}{\endcomment}
\newenvironment{history}{\comment}{\endcomment}

% For commutative diagrams we use Xy-pic
\usepackage[all]{xy}

% We use 2cell for 2-commutative diagrams.
\xyoption{2cell}
\UseAllTwocells

% We use multicol for the list of chapters between chapters
\usepackage{multicol}

% This is generally recommended for better output
\usepackage{lmodern}
\usepackage[T1]{fontenc}

% For cross-file-references
\usepackage{xr-hyper}

% Package for hypertext links:
\usepackage{hyperref}

% For any local file, say "hello.tex" you want to link to please
% use \externaldocument[hello-]{hello}
\externaldocument[introduction-]{introduction}
\externaldocument[conventions-]{conventions}
\externaldocument[sets-]{sets}
\externaldocument[categories-]{categories}
\externaldocument[topology-]{topology}
\externaldocument[sheaves-]{sheaves}
\externaldocument[sites-]{sites}
\externaldocument[stacks-]{stacks}
\externaldocument[fields-]{fields}
\externaldocument[algebra-]{algebra}
\externaldocument[brauer-]{brauer}
\externaldocument[homology-]{homology}
\externaldocument[derived-]{derived}
\externaldocument[simplicial-]{simplicial}
\externaldocument[more-algebra-]{more-algebra}
\externaldocument[smoothing-]{smoothing}
\externaldocument[modules-]{modules}
\externaldocument[sites-modules-]{sites-modules}
\externaldocument[injectives-]{injectives}
\externaldocument[cohomology-]{cohomology}
\externaldocument[sites-cohomology-]{sites-cohomology}
\externaldocument[dga-]{dga}
\externaldocument[dpa-]{dpa}
\externaldocument[sdga-]{sdga}
\externaldocument[hypercovering-]{hypercovering}
\externaldocument[schemes-]{schemes}
\externaldocument[constructions-]{constructions}
\externaldocument[properties-]{properties}
\externaldocument[morphisms-]{morphisms}
\externaldocument[coherent-]{coherent}
\externaldocument[divisors-]{divisors}
\externaldocument[limits-]{limits}
\externaldocument[varieties-]{varieties}
\externaldocument[topologies-]{topologies}
\externaldocument[descent-]{descent}
\externaldocument[perfect-]{perfect}
\externaldocument[more-morphisms-]{more-morphisms}
\externaldocument[flat-]{flat}
\externaldocument[groupoids-]{groupoids}
\externaldocument[more-groupoids-]{more-groupoids}
\externaldocument[etale-]{etale}
\externaldocument[chow-]{chow}
\externaldocument[intersection-]{intersection}
\externaldocument[pic-]{pic}
\externaldocument[weil-]{weil}
\externaldocument[adequate-]{adequate}
\externaldocument[dualizing-]{dualizing}
\externaldocument[duality-]{duality}
\externaldocument[discriminant-]{discriminant}
\externaldocument[derham-]{derham}
\externaldocument[local-cohomology-]{local-cohomology}
\externaldocument[algebraization-]{algebraization}
\externaldocument[curves-]{curves}
\externaldocument[resolve-]{resolve}
\externaldocument[models-]{models}
\externaldocument[functors-]{functors}
\externaldocument[equiv-]{equiv}
\externaldocument[pione-]{pione}
\externaldocument[etale-cohomology-]{etale-cohomology}
\externaldocument[proetale-]{proetale}
\externaldocument[relative-cycles-]{relative-cycles}
\externaldocument[more-etale-]{more-etale}
\externaldocument[trace-]{trace}
\externaldocument[crystalline-]{crystalline}
\externaldocument[spaces-]{spaces}
\externaldocument[spaces-properties-]{spaces-properties}
\externaldocument[spaces-morphisms-]{spaces-morphisms}
\externaldocument[decent-spaces-]{decent-spaces}
\externaldocument[spaces-cohomology-]{spaces-cohomology}
\externaldocument[spaces-limits-]{spaces-limits}
\externaldocument[spaces-divisors-]{spaces-divisors}
\externaldocument[spaces-over-fields-]{spaces-over-fields}
\externaldocument[spaces-topologies-]{spaces-topologies}
\externaldocument[spaces-descent-]{spaces-descent}
\externaldocument[spaces-perfect-]{spaces-perfect}
\externaldocument[spaces-more-morphisms-]{spaces-more-morphisms}
\externaldocument[spaces-flat-]{spaces-flat}
\externaldocument[spaces-groupoids-]{spaces-groupoids}
\externaldocument[spaces-more-groupoids-]{spaces-more-groupoids}
\externaldocument[bootstrap-]{bootstrap}
\externaldocument[spaces-pushouts-]{spaces-pushouts}
\externaldocument[spaces-chow-]{spaces-chow}
\externaldocument[groupoids-quotients-]{groupoids-quotients}
\externaldocument[spaces-more-cohomology-]{spaces-more-cohomology}
\externaldocument[spaces-simplicial-]{spaces-simplicial}
\externaldocument[spaces-duality-]{spaces-duality}
\externaldocument[formal-spaces-]{formal-spaces}
\externaldocument[restricted-]{restricted}
\externaldocument[spaces-resolve-]{spaces-resolve}
\externaldocument[formal-defos-]{formal-defos}
\externaldocument[defos-]{defos}
\externaldocument[cotangent-]{cotangent}
\externaldocument[examples-defos-]{examples-defos}
\externaldocument[algebraic-]{algebraic}
\externaldocument[examples-stacks-]{examples-stacks}
\externaldocument[stacks-sheaves-]{stacks-sheaves}
\externaldocument[criteria-]{criteria}
\externaldocument[artin-]{artin}
\externaldocument[quot-]{quot}
\externaldocument[stacks-properties-]{stacks-properties}
\externaldocument[stacks-morphisms-]{stacks-morphisms}
\externaldocument[stacks-limits-]{stacks-limits}
\externaldocument[stacks-cohomology-]{stacks-cohomology}
\externaldocument[stacks-perfect-]{stacks-perfect}
\externaldocument[stacks-introduction-]{stacks-introduction}
\externaldocument[stacks-more-morphisms-]{stacks-more-morphisms}
\externaldocument[stacks-geometry-]{stacks-geometry}
\externaldocument[moduli-]{moduli}
\externaldocument[moduli-curves-]{moduli-curves}
\externaldocument[examples-]{examples}
\externaldocument[exercises-]{exercises}
\externaldocument[guide-]{guide}
\externaldocument[desirables-]{desirables}
\externaldocument[coding-]{coding}
\externaldocument[obsolete-]{obsolete}
\externaldocument[fdl-]{fdl}
\externaldocument[index-]{index}

% Theorem environments.
%
\theoremstyle{plain}
\newtheorem{theorem}[subsection]{Theorem}
\newtheorem{proposition}[subsection]{Proposition}
\newtheorem{lemma}[subsection]{Lemma}

\theoremstyle{definition}
\newtheorem{definition}[subsection]{Definition}
\newtheorem{example}[subsection]{Example}
\newtheorem{exercise}[subsection]{Exercise}
\newtheorem{situation}[subsection]{Situation}

\theoremstyle{remark}
\newtheorem{remark}[subsection]{Remark}
\newtheorem{remarks}[subsection]{Remarks}

\numberwithin{equation}{subsection}

% Macros
%
\def\lim{\mathop{\mathrm{lim}}\nolimits}
\def\colim{\mathop{\mathrm{colim}}\nolimits}
\def\Spec{\mathop{\mathrm{Spec}}}
\def\Hom{\mathop{\mathrm{Hom}}\nolimits}
\def\Ext{\mathop{\mathrm{Ext}}\nolimits}
\def\SheafHom{\mathop{\mathcal{H}\!\mathit{om}}\nolimits}
\def\SheafExt{\mathop{\mathcal{E}\!\mathit{xt}}\nolimits}
\def\Sch{\mathit{Sch}}
\def\Mor{\mathop{\mathrm{Mor}}\nolimits}
\def\Ob{\mathop{\mathrm{Ob}}\nolimits}
\def\Sh{\mathop{\mathit{Sh}}\nolimits}
\def\NL{\mathop{N\!L}\nolimits}
\def\CH{\mathop{\mathrm{CH}}\nolimits}
\def\proetale{{pro\text{-}\acute{e}tale}}
\def\etale{{\acute{e}tale}}
\def\QCoh{\mathit{QCoh}}
\def\Ker{\mathop{\mathrm{Ker}}}
\def\Im{\mathop{\mathrm{Im}}}
\def\Coker{\mathop{\mathrm{Coker}}}
\def\Coim{\mathop{\mathrm{Coim}}}

% Boxtimes
%
\DeclareMathSymbol{\boxtimes}{\mathbin}{AMSa}{"02}

%
% Macros for moduli stacks/spaces
%
\def\QCohstack{\mathcal{QC}\!\mathit{oh}}
\def\Cohstack{\mathcal{C}\!\mathit{oh}}
\def\Spacesstack{\mathcal{S}\!\mathit{paces}}
\def\Quotfunctor{\mathrm{Quot}}
\def\Hilbfunctor{\mathrm{Hilb}}
\def\Curvesstack{\mathcal{C}\!\mathit{urves}}
\def\Polarizedstack{\mathcal{P}\!\mathit{olarized}}
\def\Complexesstack{\mathcal{C}\!\mathit{omplexes}}
% \Pic is the operator that assigns to X its picard group, usage \Pic(X)
% \Picardstack_{X/B} denotes the Picard stack of X over B
% \Picardfunctor_{X/B} denotes the Picard functor of X over B
\def\Pic{\mathop{\mathrm{Pic}}\nolimits}
\def\Picardstack{\mathcal{P}\!\mathit{ic}}
\def\Picardfunctor{\mathrm{Pic}}
\def\Deformationcategory{\mathcal{D}\!\mathit{ef}}
\usepackage{cleveref}
\crefname{originallemma}{original lemma}{original lemmas}
\Crefname{originallemma}{Original lemma}{Original lemmas}
\makeatletter
\newlabel{algebra-lemma-universally-exact-split}{{058L}{source}{}{sourcepointer-1}{}}
\newlabel{algebra-lemma-universally-exact-split@cref}{{[originallemma][0][]058L}{[1][0][]source}{}{}{}}
\makeatother
\crefname{originaltheorem}{original theorem}{original theorems}
\Crefname{originaltheorem}{Original theorem}{Original theorems}
\makeatletter
\newlabel{algebra-theorem-universally-exact-criteria}{{058K}{source}{}{sourcepointer-2}{}}
\newlabel{algebra-theorem-universally-exact-criteria@cref}{{[originaltheorem][0][]058K}{[1][0][]source}{}{}{}}
\makeatother
\crefname{originallemma}{original lemma}{original lemmas}
\Crefname{originallemma}{Original lemma}{Original lemmas}
\makeatletter
\newlabel{derived-lemma-compute-ext-resolutions}{{06XR}{source}{}{sourcepointer-3}{}}
\newlabel{derived-lemma-compute-ext-resolutions@cref}{{[originallemma][0][]06XR}{[1][0][]source}{}{}{}}
\makeatother
\crefname{originallemma}{original lemma}{original lemmas}
\Crefname{originallemma}{Original lemma}{Original lemmas}
\makeatletter
\newlabel{derived-lemma-yoneda-extension}{{06XU}{source}{}{sourcepointer-4}{}}
\newlabel{derived-lemma-yoneda-extension@cref}{{[originallemma][0][]06XU}{[1][0][]source}{}{}{}}
\makeatother
\crefname{originalreference}{original reference}{original references}
\Crefname{originalreference}{Original reference}{Original references}
\makeatletter
\newlabel{item-mult-map}{{06Z9}{source}{}{sourcepointer-5}{}}
\newlabel{item-mult-map@cref}{{[originalreference][0][]06Z9}{[1][0][]source}{}{}{}}
\makeatother
\crefname{originalreference}{original reference}{original references}
\Crefname{originalreference}{Original reference}{Original references}
\makeatletter
\newlabel{item-mult-map-linear}{{06ZC}{source}{}{sourcepointer-6}{}}
\newlabel{item-mult-map-linear@cref}{{[originalreference][0][]06ZC}{[1][0][]source}{}{}{}}
\makeatother
\crefname{originalreference}{original reference}{original references}
\Crefname{originalreference}{Original reference}{Original references}
\makeatletter
\newlabel{item-tangent-right-exact}{{06ZD}{source}{}{sourcepointer-7}{}}
\newlabel{item-tangent-right-exact@cref}{{[originalreference][0][]06ZD}{[1][0][]source}{}{}{}}
\makeatother
\crefname{originallemma}{original lemma}{original lemmas}
\Crefname{originallemma}{Original lemma}{Original lemmas}
\makeatletter
\newlabel{lemma-enough-injectives}{{06Z7}{source}{}{sourcepointer-8}{}}
\newlabel{lemma-enough-injectives@cref}{{[originallemma][0][]06Z7}{[1][0][]source}{}{}{}}
\makeatother
\crefname{originallemma}{original lemma}{original lemmas}
\Crefname{originallemma}{Original lemma}{Original lemmas}
\makeatletter
\newlabel{lemma-exactness-implies}{{06ZB}{source}{}{sourcepointer-9}{}}
\newlabel{lemma-exactness-implies@cref}{{[originallemma][0][]06ZB}{[1][0][]source}{}{}{}}
\makeatother
\crefname{originallemma}{original lemma}{original lemmas}
\Crefname{originallemma}{Original lemma}{Original lemmas}
\makeatletter
\newlabel{lemma-exactness-permanence}{{06ZF}{source}{}{sourcepointer-10}{}}
\newlabel{lemma-exactness-permanence@cref}{{[originallemma][0][]06ZF}{[1][0][]source}{}{}{}}
\makeatother
\crefname{originallemma}{original lemma}{original lemmas}
\Crefname{originallemma}{Original lemma}{Original lemmas}
\makeatletter
\newlabel{lemma-tangent-functor}{{06Z8}{source}{}{sourcepointer-11}{}}
\newlabel{lemma-tangent-functor@cref}{{[originallemma][0][]06Z8}{[1][0][]source}{}{}{}}
\makeatother
\crefname{originallemma}{original lemma}{original lemmas}
\Crefname{originallemma}{Original lemma}{Original lemmas}
\makeatletter
\newlabel{lemma-tangent-injective}{{06ZA}{source}{}{sourcepointer-12}{}}
\newlabel{lemma-tangent-injective@cref}{{[originallemma][0][]06ZA}{[1][0][]source}{}{}{}}
\makeatother
\title{The Ext groups out of a finitely presented functor on modules vanish in positive degrees}
\author{The Stacks Project Authors}
\date{stacks-project:adequate.tex:local-snapshot-20261012; GNU FDL 1.2 or later}
\begin{document}
\maketitle
\noindent\textit{Editorial note.} All delivered mathematical text is the author's own, copied from the pinned Stacks Project chapter file adequate.tex; no statement or proof step was rewritten, reordered or removed.\par\medskip
\noindent\textit{Editorial note.} The counted spans carry navigation stubs for the references leaving them; no bibliography citation occurs in the spans.\par\medskip
\noindent\textit{Editorial note.} The preamble span points at the chapter's own preamble.tex shipped with the frozen source, used verbatim as staged.\par\medskip
\noindent\textit{Allowed standard prerequisites (editorial).}\begin{itemize}
\item The category of module-valued functors on Alg\_A has enough injectives. Same chapter, Lemma lemma-enough-injectives, invoked at proof lines 1434-1436. Provides the injective resolution used to compute the Ext groups.
\item Ext computed from resolutions, and the Yoneda description of Ext classes. Derived Categories, Lemmas derived-lemma-compute-ext-resolutions and derived-lemma-yoneda-extension, invoked at proof lines 1437-1449. Turns an Ext class into a morphism into the resolution and back into a Yoneda extension.
\item Property (*) and its permanence for kernels, images and fibre products of module-valued functors. Same chapter, Lemmas lemma-tangent-injective, lemma-exactness-permanence and lemma-exactness-implies, invoked at proof lines 1456-1496. Makes every functor in the extension satisfy (*) and gives the B-linearity of the comparison maps.
\item The tangent functor TF and its multiplication map. Same chapter, Lemma lemma-tangent-functor (item item-mult-map), invoked at proof lines 1489-1496. Supplies the canonical maps F\_j(B) -> G\_j(B) used in the comparison.
\end{itemize}\medskip
\ifdefined\hypertarget\hypertarget{sourcepointer-1}{}\fi
\noindent\textit{Source reference.} Original lemma 058L: sources/stacks-chow/tags.txt, source lines 6784--6784. Pointer outside the counted span, invoked once for context; the load-bearing argument of the counted statement is the author's proof reproduced in full in the proof span.\par\medskip
\ifdefined\hypertarget\hypertarget{sourcepointer-2}{}\fi
\noindent\textit{Source reference.} Original theorem 058K: sources/stacks-chow/tags.txt, source lines 6783--6783. Pointer outside the counted span, invoked once for context; the load-bearing argument of the counted statement is the author's proof reproduced in full in the proof span.\par\medskip
\ifdefined\hypertarget\hypertarget{sourcepointer-3}{}\fi
\noindent\textit{Source reference.} Original lemma 06XR: sources/stacks-chow/tags.txt, source lines 8860--8860. Pointer outside the counted span, invoked once for context; the load-bearing argument of the counted statement is the author's proof reproduced in full in the proof span.\par\medskip
\ifdefined\hypertarget\hypertarget{sourcepointer-4}{}\fi
\noindent\textit{Source reference.} Original lemma 06XU: sources/stacks-chow/tags.txt, source lines 8863--8863. Pointer outside the counted span, invoked once for context; the load-bearing argument of the counted statement is the author's proof reproduced in full in the proof span.\par\medskip
\ifdefined\hypertarget\hypertarget{sourcepointer-5}{}\fi
\noindent\textit{Source reference.} Original reference 06Z9: sources/stacks-adequate/pinned-original-source.tex, source lines 1152--1152. Pointer outside the counted span, invoked once for context; the load-bearing argument of the counted statement is the author's proof reproduced in full in the proof span.\par\medskip
\ifdefined\hypertarget\hypertarget{sourcepointer-6}{}\fi
\noindent\textit{Source reference.} Original reference 06ZC: sources/stacks-adequate/pinned-original-source.tex, source lines 1282--1282. Pointer outside the counted span, invoked once for context; the load-bearing argument of the counted statement is the author's proof reproduced in full in the proof span.\par\medskip
\ifdefined\hypertarget\hypertarget{sourcepointer-7}{}\fi
\noindent\textit{Source reference.} Original reference 06ZD: sources/stacks-adequate/pinned-original-source.tex, source lines 1287--1287. Pointer outside the counted span, invoked once for context; the load-bearing argument of the counted statement is the author's proof reproduced in full in the proof span.\par\medskip
\ifdefined\hypertarget\hypertarget{sourcepointer-8}{}\fi
\noindent\textit{Source reference.} Original lemma 06Z7: sources/stacks-adequate/pinned-original-source.tex, source lines 1086--1086. Pointer outside the counted span, invoked once for context; the load-bearing argument of the counted statement is the author's proof reproduced in full in the proof span.\par\medskip
\ifdefined\hypertarget\hypertarget{sourcepointer-9}{}\fi
\noindent\textit{Source reference.} Original lemma 06ZB: sources/stacks-adequate/pinned-original-source.tex, source lines 1271--1271. Pointer outside the counted span, invoked once for context; the load-bearing argument of the counted statement is the author's proof reproduced in full in the proof span.\par\medskip
\ifdefined\hypertarget\hypertarget{sourcepointer-10}{}\fi
\noindent\textit{Source reference.} Original lemma 06ZF: sources/stacks-adequate/pinned-original-source.tex, source lines 1359--1359. Pointer outside the counted span, invoked once for context; the load-bearing argument of the counted statement is the author's proof reproduced in full in the proof span.\par\medskip
\ifdefined\hypertarget\hypertarget{sourcepointer-11}{}\fi
\noindent\textit{Source reference.} Original lemma 06Z8: sources/stacks-adequate/pinned-original-source.tex, source lines 1138--1138. Pointer outside the counted span, invoked once for context; the load-bearing argument of the counted statement is the author's proof reproduced in full in the proof span.\par\medskip
\ifdefined\hypertarget\hypertarget{sourcepointer-12}{}\fi
\noindent\textit{Source reference.} Original lemma 06ZA: sources/stacks-adequate/pinned-original-source.tex, source lines 1195--1195. Pointer outside the counted span, invoked once for context; the load-bearing argument of the counted statement is the author's proof reproduced in full in the proof span.\par\medskip
\begin{lemma}
\label{lemma-ext-group-zero-key}
Let $A$ be a ring. Let $M$, $P$ be $A$-modules with $P$ of finite
presentation. Then
$\Ext^i_\mathcal{P}(\underline{P}, \underline{M}) = 0$
for $i > 0$ where $\mathcal{P}$ is the category of module-valued
functors on $\textit{Alg}_A$.
\end{lemma}

\begin{proof}
Choose an injective resolution $\underline{M} \to I^\bullet$ in
$\mathcal{P}$, see
Lemma \ref{lemma-enough-injectives}.
By
Derived Categories, Lemma \ref{derived-lemma-compute-ext-resolutions}
any element of $\Ext^i_\mathcal{P}(\underline{P}, \underline{M})$
comes from a morphism $\varphi : \underline{P} \to I^i$ with
$d^i \circ \varphi = 0$. We will prove that the
Yoneda extension
$$
E : 0 \to \underline{M} \to I^0 \to \ldots \to
I^{i - 1} \times_{\Ker(d^i)} \underline{P} \to \underline{P} \to 0
$$
of $\underline{P}$ by $\underline{M}$
associated to $\varphi$ is trivial, which will prove the lemma by
Derived Categories, Lemma \ref{derived-lemma-yoneda-extension}.

\medskip\noindent
For $F$ a module-valued functor on $\textit{Alg}_A$
say $(*)$ holds if for all $B \in \Ob(\textit{Alg}_A)$ the
functor $TF(B, -)$ on $B$-modules transforms a short exact sequence
of $B$-modules into a right exact sequence.
Recall that the module-valued functors $\underline{M}, I^n, \underline{P}$
each have property $(*)$, see
Lemma \ref{lemma-tangent-injective}
and the remarks preceding it.
By splitting $0 \to \underline{M} \to I^\bullet$ into short
exact sequences we find that each of the functors
$\Im(d^{n - 1}) = \Ker(d^n) \subset I^n$ has property $(*)$ by
Lemma \ref{lemma-exactness-permanence}
and also that $I^{i - 1} \times_{\Ker(d^i)} \underline{P}$ has property
$(*)$.

\medskip\noindent
Thus we may assume the Yoneda extension is given as
$$
E : 0 \to \underline{M} \to F_{i - 1} \to \ldots \to
F_0 \to \underline{P} \to 0
$$
where each of the module-valued functors $F_j$ has property $(*)$.
Set $G_j(B) = TF_j(B, B)$ viewed as a $B$-module via the {\it second}
$B$-module structure defined in
Lemma \ref{lemma-exactness-implies}.
Since $TF_j$ is a functor on pairs we see that $G_j$ is a module-valued
functor on $\textit{Alg}_A$. Moreover, since $E$ is an exact sequence
the sequence $G_{j + 1} \to G_j \to G_{j - 1}$ is exact (see remark
preceding
Lemma \ref{lemma-exactness-permanence}).
Observe that $T\underline{M}(B, B) = M \otimes_A B = \underline{M}(B)$
and that the two $B$-module structures agree on this.
Thus we obtain a Yoneda extension
$$
E' :  0 \to \underline{M} \to G_{i - 1} \to \ldots \to
G_0 \to \underline{P} \to 0
$$
Moreover, the canonical maps
$$
F_j(B) = B \otimes_B F_j(B) \longrightarrow TF_j(B, B) = G_j(B)
$$
of
Lemma \ref{lemma-tangent-functor} (\ref{item-mult-map})
are $B$-linear by
Lemma \ref{lemma-exactness-implies} (\ref{item-mult-map-linear})
and functorial in $B$. Hence a map
$$
\xymatrix{
0 \ar[r] &
\underline{M} \ar[r] \ar[d]^1 &
F_{i - 1} \ar[r] \ar[d] &
\ldots \ar[r] &
F_0 \ar[r] \ar[d] &
\underline{P} \ar[r] \ar[d]^1 & 0 \\
0 \ar[r] &
\underline{M} \ar[r] &
G_{i - 1} \ar[r] &
\ldots \ar[r] &
G_0 \ar[r] &
\underline{P} \ar[r] & 0
}
$$
of Yoneda extensions. In particular we see that $E$ and $E'$ have the
same class in $\Ext^i_\mathcal{P}(\underline{P}, \underline{M})$
by the lemma on Yoneda Exts mentioned above. Finally, let $N$ be a
$A$-module of finite presentation. Then we see that
$$
0 \to T\underline{M}(A, N) \to TF_{i - 1}(A, N) \to \ldots \to
TF_0(A, N) \to T\underline{P}(A, N) \to 0
$$
is exact. By
Lemma \ref{lemma-exactness-implies} (\ref{item-tangent-right-exact})
with $B = A$ this translates into the exactness of the sequence of
$A$-modules
$$
0 \to M \otimes_A N \to G_{i - 1}(A) \otimes_A N \to \ldots \to
G_0(A) \otimes_A N \to P \otimes_A N \to 0
$$
Hence the sequence of $A$-modules
$0 \to M \to G_{i - 1}(A) \to \ldots \to G_0(A) \to P \to 0$
is universally exact, in the sense that it remains exact on tensoring
with any finitely presented $A$-module $N$. Let
$K = \Ker(G_0(A) \to P)$ so that we have exact sequences
$$
0 \to K \to G_0(A) \to P \to 0
\quad\text{and}\quad
G_2(A) \to G_1(A) \to K \to 0
$$
Tensoring the second sequence with $N$ we obtain that
$K \otimes_A N = \Coker(G_2(A) \otimes_A N \to G_1(A) \otimes_A N)$.
Exactness of $G_2(A) \otimes_A N \to G_1(A) \otimes_A N \to G_0(A) \otimes_A N$
then implies that $K \otimes_A N \to G_0(A) \otimes_A N$ is injective.
By
Algebra, Theorem \ref{algebra-theorem-universally-exact-criteria}
this means that the $A$-module extension $0 \to K \to G_0(A) \to P \to 0$
is exact, and because $P$ is assumed of finite presentation this means
the sequence is split, see
Algebra, Lemma \ref{algebra-lemma-universally-exact-split}.
Any splitting $P \to G_0(A)$ defines a map $\underline{P} \to G_0$
which splits the surjection $G_0 \to \underline{P}$. Thus the
Yoneda extension $E'$ is equivalent to the trivial Yoneda extension
and we win.
\end{proof}

\end{document}
